% theme: similarity_ratios_and_volume_ratios
% 叩き台。解答は未生成。レビュー用図を挿入済み。表3大問、裏4大問。
% 体積の比の逆算は既知の整数比の3乗で処理できる数値に限定する。
\MajorQuestionLesson{相似比と体積比}
\SetRowSpace{3mm}

%表面は錐体を底面に平行に切る問題


\TypeQuestion{}{similarity_volume_ratio_pyramid_cut}
\FigProblemBlock{48mm}{fig/lesson_11_pyramid_cut}{
\Text{高さが $12\mathrm{cm}$、底面積が $81\mathrm{cm}^2$ の正四角錐を、底面に平行な平面で切る。頂点から切断面までの距離は $8\mathrm{cm}$ である。}
}{
\QQ{切断面の面積を求めなさい。}{}
\QQ{頂点を含む小さい正四角錐の体積を求めなさい。}{}
\QQ{切断後の2つの立体のうち、頂点を含まない方の体積を求めなさい。}{}
}

%１度切断ではなく、２〜３回切断にしてください。
\TypeQuestion{}{similarity_volume_ratio_inverse_cut}
\FigProblemBlock{48mm}{fig/lesson_11_inverse_cut}{
\Text{高さが $18\mathrm{cm}$、体積が $648\mathrm{cm}^3$ の正四角錐を、底面に平行な平面で切った。頂点を含まない方の立体の体積は $456\mathrm{cm}^3$ であった。}
}{
\QQ{頂点を含む小さい正四角錐の体積を求めなさい。}{}
\QQ{頂点から切断面までの距離を求めなさい。}{}
\QQ{切断面からもとの底面までの距離を求めなさい。}{}


\TypeQuestion{}{similarity_volume_ratio_water_depth}
\FigProblemBlock{48mm}{fig/lesson_11_water}{
\Text{頂点を下にして置いた円錐形の容器がある。容器の深さは $12\mathrm{cm}$、満水のときの水の体積は $960\mathrm{cm}^3$ である。水面は容器の上端の面と平行である。容器の厚さは考えないものとする。}
}{
\QQ{水深が $6\mathrm{cm}$ のとき、水の体積を求めなさい。}{}
\QQ{$405\mathrm{cm}^3$ の水を入れたとき、水深を求めなさい。}{}
\QQ{深さが $12\mathrm{cm}$、満水のときの水の体積が $960\mathrm{cm}^3$ の円柱形の容器を、底面が水平になるように置く。2つの容器の水深がともに $6\mathrm{cm}$ のとき、円柱形の容器に入る水の体積は、円錐形の容器に入る水の体積の何倍か求めなさい。}{}
}
}
\AnswerColumnBreak
\LessonQuestionBreak


%裏：送った図のような感じ。
%大問＋小問構成
%立方体や直方体を切断し、角錐台を作る問題


